\section{Sensitivity to the data window}
\label{sec:window-comparison}
The current grid replaces the earlier fitting selection; it does not
discard the older observations. We retain seven stages in the CSV:
\begin{itemize}
\item A (\texttt{original}): $N=16,18,20,22,24$, $\ell=2,\ldots,6$,
$m=13$ after the cutoff.
\item B (\texttt{ell1}): those sizes with $\ell=1,\ldots,6$, $m=18$.
\item C (\texttt{extended}): B plus $N=28,32,48,64,96,128$,
$\ell=1,\ldots,6$, $m=53$.
\item D (\texttt{large\_intervals}): $N=32,48,64,82,96,112,128$,
$\ell=4,\ldots,10$, with $x\leq0.20$ and $m=44$.
\item E (\texttt{larger\_chains\_wide}): all 15 sizes in
Eq.~\eqref{eq:large-grid}, still using $x\leq0.20$, with $m=100$.
\item F (\texttt{larger\_chains}): the same computed grid, with the
previous fitting cutoff $x\leq0.12$ and $m=89$.
\item G (\texttt{narrow\_008}): the same computed grid, with the
current fitting cutoff $x\leq0.08$ and $m=77$.
\end{itemize}
All stages have $p=1/2$, the same fixed powers, and unweighted least
squares in $D$. Stages A--E use $x\leq0.20$; F uses $x\leq0.12$;
G uses $x\leq0.08$.
All $F_1$--$F_4$ fit and analytical CFT series
summaries are stored in \texttt{stage\_comparison.csv}.
Tables~\ref{tab:ising-window-comparison} and \ref{tab:boson-window-comparison}
highlight $F_3$ for E, F, G, comparing three cutoffs on the same
computed grid. Their validation samples differ, so the
errors compare sensitivity to the window, not performance on an identical
test set. Superscripts on coefficients label the stage, not an exponent.

\begin{table}[htbp]\centering\scriptsize
\begin{tabular}{lrrrrrr}\toprule
Direction & $c_\alpha^{(E)}$ & $c_\alpha^{(F)}$ & $c_\alpha^{(G)}$ & CV RMSE E & CV RMSE F & CV RMSE G\\\midrule
$I\mid \sigma$ & 0.020849 & 0.020847 & 0.020856 & $1.92\!\times\!10^{-6}$ & $6.20\!\times\!10^{-7}$ & $3.56\!\times\!10^{-7}$\\
$\sigma\mid I$ & 0.020845 & 0.020845 & 0.020852 & $9.98\!\times\!10^{-7}$ & $3.55\!\times\!10^{-7}$ & $2.23\!\times\!10^{-7}$\\
$I\mid \varepsilon$ & 0.12937 & 0.13198 & 0.13058 & $1.62\!\times\!10^{-5}$ & $2.82\!\times\!10^{-6}$ & $1.14\!\times\!10^{-6}$\\
$\varepsilon\mid I$ & 0.12422 & 0.13095 & 0.13044 & $1.36\!\times\!10^{-5}$ & $2.45\!\times\!10^{-6}$ & $1.03\!\times\!10^{-6}$\\
$\sigma\mid \varepsilon$ & 0.021364 & 0.020891 & 0.020772 & $3.23\!\times\!10^{-5}$ & $6.79\!\times\!10^{-6}$ & $3.08\!\times\!10^{-6}$\\
$\varepsilon\mid \sigma$ & 0.021692 & 0.020978 & 0.020785 & $5.99\!\times\!10^{-5}$ & $9.25\!\times\!10^{-6}$ & $3.63\!\times\!10^{-6}$\\
\bottomrule\end{tabular}

\caption{Ising $F_3$ fits for stages E, F, G, with three free coefficients
at powers $\alpha,\alpha+2,\alpha+4$ and fixed
$\alpha=4$ for $I\leftrightarrow\eps$, two otherwise. The geometries,
sample counts and stage-dependent cutoffs are defined in
Appendix~\ref{sec:window-comparison}. States use the
critical periodic Hamiltonian $H=-\sum XX-\sum Z$ and exact covariance
reductions. CV RMSE is the complete-size holdout error in
Eq.~\eqref{eq:cv}. Coefficients multiply powers of $\pi x$.}
\label{tab:ising-window-comparison}
\end{table}
\begin{table}[htbp]\centering\scriptsize
\begin{tabular}{lrrrrrr}\toprule
Direction & $c_\alpha^{(E)}$ & $c_\alpha^{(F)}$ & $c_\alpha^{(G)}$ & CV RMSE E & CV RMSE F & CV RMSE G\\\midrule
$00\mid ss$ & 0.078378 & 0.078307 & 0.076917 & $1.76\!\times\!10^{-4}$ & $7.28\!\times\!10^{-5}$ & $4.81\!\times\!10^{-5}$\\
$ss\mid 00$ & 0.07837 & 0.078306 & 0.076919 & $1.74\!\times\!10^{-4}$ & $7.25\!\times\!10^{-5}$ & $4.80\!\times\!10^{-5}$\\
$00\mid 3s,s$ & 0.39905 & 0.39904 & 0.39364 & $6.10\!\times\!10^{-4}$ & $2.78\!\times\!10^{-4}$ & $1.86\!\times\!10^{-4}$\\
$3s,s\mid 00$ & 0.39899 & 0.39897 & 0.39376 & $5.26\!\times\!10^{-4}$ & $2.61\!\times\!10^{-4}$ & $1.79\!\times\!10^{-4}$\\
$ss\mid 3s,s$ & 0.16424 & 0.16422 & 0.1628 & $1.77\!\times\!10^{-4}$ & $7.31\!\times\!10^{-5}$ & $4.71\!\times\!10^{-5}$\\
$3s,s\mid ss$ & 0.16424 & 0.16415 & 0.16288 & $1.07\!\times\!10^{-4}$ & $6.18\!\times\!10^{-5}$ & $4.29\!\times\!10^{-5}$\\
\bottomrule\end{tabular}

\caption{Boson $F_3$ fits for the same stages E, F, G and objective, with
fixed $\alpha=2$, three free coefficients, charge codes $(0,0),(1,1),(3,1)$,
$q=2$, $z_j=e^{2\pi i j/N}$, $w_1=0$, $W=10^8$, and normalized
vertex/Jastrow states. All stages use $p=1/2$ and natural logarithms;
CV RMSE has the definition in Eq.~\eqref{eq:cv}.}
\label{tab:boson-window-comparison}
\end{table}

For Ising $\sigma\mid\eps$, $c_2$ changes from $0.021364$ (E) to
$0.020891$ (F) and $0.020772$ (G), compared with $1/48$ from CFT.
For the boson $ss\mid3s,s$, the values are $0.164240$, $0.164217$,
and $0.162805$, compared with $1/6$. The tighter cutoff therefore does
not uniformly improve agreement with the theoretical leading coefficient.
All six boson $F_3$ fits in G have worse size-holdout errors than $F_1$;
all six Ising directions retain a positive CV gain. Moreover,
least squares in $D$ gives small absolute-$D$ points less influence than
a relative-error objective would, and the fitted coefficients remain
finite-window estimates rather than controlled continuum extrapolations.
